\chapter{Finestra Dominio Temporale (TemporalWindow)}

\section{Scopo}
Definisce l'intervallo temporale di simulazione (inizio/fine) e i parametri derivati.

\section{Funzioni principali}
\begin{itemize}[leftmargin=1.5em]
  \item Selezione data e ora inizio simulazione.
  \item Selezione data e ora fine simulazione.
  \item Calcolo e verifica della durata complessiva.
  \item Salvataggio in \texttt{temp\_config/temporal\_config.json}.
\end{itemize}

\section{Sottofinestre}
\subsection{Calendario (open\_calendar)}
La finestra apre un calendario in \texttt{tk.Toplevel} per semplificare la scelta delle date.
Il dialog viene usato sia per data di inizio sia per data di fine.

\section{Output e integrazione}
Il file temporale viene letto dai moduli di generazione input CALMET/CALPUFF/CALPOST e dalle operazioni sul farm.
